\section{Introduction and Threat Model}\label{sec:intro}
A client of an LLM API names a model, pays that model's per-token price, and gets back text along
with a response that repeats the model's name. The repeated name proves nothing. A gateway,
reseller or router between the client and the provider can serve a smaller model, a quantised
copy, or a mix of the two, and still return the name the client asked for. The incentive is large.
For the token mix used in this paper the substitute costs 11.5\% of the genuine model, so a gateway
that swaps one request in ten saves about \$88 per \$1,000 of billed traffic.

The client pays for model $M$, and the only record of what actually ran is the gateway's own. A
provider could attest to the computation in trusted hardware~\cite{cai2025}, but gateways do not
offer this today, and the client has no way to inspect them. What the client can do is test the
responses. A black-box \emph{auditor} sends probe requests and decides, from the text, timing and
metadata that come back, whether the advertised model answered.

At least eight such auditors have appeared since April
2025~\cite{cai2025,zhu2025rut,lin2026gatescope,fang2026kbf,bruckner2026,iris2026,difr2025}. Each
was evaluated on its own testbed and against its own adversary, usually with thresholds chosen on
the same data it reports. A frozen-threshold replication~\cite{holdout2026} shows how fragile this
is: a method with perfect development accuracy fell to 50\% sensitivity on fresh model pairs, and
half of those pairs could not be scored at all. The auditors have never been compared on a common
testbed against a common adversary, and none has been tested against a gateway that expects to be
audited.

\textbf{Threat model.} The client sends requests over HTTPS to a gateway that advertises model $M$.
The gateway controls routing and response metadata. For each request it may forward to $M$, to a
cheaper substitute $M'$, to another provider of $M$, or to a cache. It may rewrite the reported
model name, the \texttt{usage} token counts and \texttt{system\_fingerprint}, and it may add delay.
Because it reads every request, it can classify requests by their shape. It does not know a given
auditor's probes or thresholds, but it knows the published methods and the public benchmarks. Its
goal is to save as much as possible while keeping the probability of detection low. The client
sees only what crosses the wire: response text, client-side latency, the \texttt{usage} block, the
fingerprint, and logprobs where the endpoint exposes them. Timing is a side channel the client can
measure without the gateway's cooperation; content is the only other signal. Honest variation, such
as load balancing across legitimate providers of $M$, is permitted and must not be flagged.

\textbf{Contributions.}
\begin{itemize}
\item \textbf{A shared adversarial testbed.} SHIM is an OpenAI-compatible gateway with twelve
  adversary arms and a hidden ledger of which backend really answered. Four of the arms
  (probe-aware evasion, laundering, canary-awareness and benign routing) appear in no prior
  evaluation.
\item \textbf{A sealed comparison of six auditors.} We reimplement five published methods and add an
  e-value fusion, all evaluated under thresholds hashed and tagged before evaluation. An auditor
  that cannot run reports \emph{uninformative}, never \emph{pass}.
\item \textbf{Deployability on real endpoints.} On 11 free-tier endpoints, a logprob-based auditor
  runs on none and the one-token method as published on 4, because hidden reasoning tokens use up
  its token budget.
\item \textbf{Three negative results.} Probe-aware evasion keeps every auditor at chance.
  False-positive budgets, not audit cost, set a dilution floor of $\epsilon^* \approx 0.10$. Timing
  auditors cannot tell fraud from sanctioned routing, which splits ``substitution'' into identity,
  disclosure and billing violations.
\end{itemize}
